% ============================================================
%  Non-Convex Polygon Packing via Simulated Annealing
%  and Bracketing-Bisection Search
%
%  Joel Solis, University of Arizona
%  Applied Mathematics and Numerical Methods, 2026
% ============================================================

\documentclass[10pt,twocolumn]{article}

% ---- Typography ------------------------------------------------
\usepackage{mathptmx}
\usepackage[T1]{fontenc}
\usepackage{microtype}
\linespread{0.96}

% ---- Page layout -----------------------------------------------
\usepackage[
  left=1.65cm, right=1.65cm,
  top=2.6cm,   bottom=2.4cm,
  headheight=13pt, headsep=10pt
]{geometry}
\setlength{\columnsep}{18pt}
\setlength{\columnseprule}{0.35pt}

% ---- Running headers -------------------------------------------
\usepackage{fancyhdr}
\pagestyle{fancy}
\fancyhf{}
\fancyhead[LO]{\small\itshape Non-Convex Polygon Packing via Simulated Annealing}
\fancyhead[RE]{\small\itshape J. Solis}
\fancyhead[RO,LE]{\small\thepage}
\renewcommand{\headrulewidth}{0.4pt}

\fancypagestyle{firstpage}{%
  \fancyhf{}
  \fancyhead[L]{\small\itshape
    Applied Mathematics and Numerical Methods --- University of Arizona, 2026}
  \fancyhead[R]{\small\thepage}
  \renewcommand{\headrulewidth}{0.4pt}%
}

% ---- Section headings ------------------------------------------
\usepackage{titlesec}
\titleformat{\section}[block]
  {\normalfont\bfseries\small\centering}{\thesection.}{0.4em}{\MakeUppercase}
\titleformat{\subsection}[block]
  {\normalfont\itshape\small}{\thesubsection.}{0.4em}{}
\titlespacing*{\section}   {0pt}{7pt plus 2pt minus 1pt}{3pt plus 1pt}
\titlespacing*{\subsection}{0pt}{5pt plus 1pt minus 1pt}{2pt plus 1pt}

% ---- Math and floats -------------------------------------------
\usepackage{amsmath,amssymb}
\usepackage[draft]{graphicx}
\usepackage{booktabs,array}
\usepackage[hidelinks]{hyperref}
\usepackage[font=small,labelfont=bf,labelsep=period,
            justification=justified]{caption}
\setlength{\abovecaptionskip}{4pt}
\setlength{\belowcaptionskip}{0pt}
\usepackage{enumitem}
\setlist{noitemsep,topsep=1pt,leftmargin=*}

% ---- Abstract --------------------------------------------------
\usepackage{abstract}
\renewcommand{\abstractname}{}
\renewcommand{\absnamepos}{empty}
\setlength{\absleftindent}{0pt}
\setlength{\absrightindent}{0pt}

% ---- Spacing ---------------------------------------------------
\setlength{\parskip}{2pt}
\setlength{\parindent}{11pt}
\setlength{\abovedisplayskip}{5pt plus 2pt minus 2pt}
\setlength{\belowdisplayskip}{5pt plus 2pt minus 2pt}

% ================================================================
\title{%
  \Large\bfseries
  Non-Convex Polygon Packing via Simulated Annealing\\[4pt]
  and Bracketing-Bisection Search%
}

\author{%
  Joel Solis%
  \thanks{Department of Mathematics, University of Arizona.
    \texttt{jsolis@math.arizona.edu}.
    Supported by allocation \texttt{ece569} on the
    University of Arizona HPC cluster.}%
}
\date{April 2026}

% ================================================================
\begin{document}
\maketitle
\thispagestyle{firstpage}

% ----------------------------------------------------------------
\begin{abstract}
\vspace{-2pt}
\noindent\rule{\linewidth}{0.4pt}\vspace{4pt}

\noindent{\bfseries\small Abstract.}\enspace\small
Take a fixed non-convex polygon. The problem, featured in a 2024
Kaggle competition, is to find the minimum side length $L^*$ of a
square container that can hold $N$ congruent copies without overlap,
each free to translate and rotate. Answering this exactly is NP-hard.
We describe a practical solver combining time-budgeted Simulated
Annealing with an outer bisection loop that narrows the bracket
containing $L^*$ to a relative precision of $10^{-3}$, followed by a
feasibility-polish stage. Collision detection uses the Separating Axis
Theorem (SAT), chosen because it returns penetration depth, converting
a combinatorial feasibility question into a smooth, navigable energy
function. We deployed embarrassingly parallel jobs on a five-node
HPC cluster for $N \in \{5, 10, 25, 50, 100, 200\}$, with wall times
from 3~minutes to 8~hours. Packing densities range from
$\eta \approx 0.58$ at small $N$ to $\eta \approx 0.48$ at $N = 200$,
with a $1/\!\sqrt{N}$ boundary-effect law supported by a geometric
argument.

\vspace{5pt}
\noindent{\bfseries\small Keywords.}\enspace\small
polygon packing; simulated annealing; separating axis theorem;
bisection search; embarrassingly parallel HPC; packing density scaling.

\vspace{4pt}
\noindent\rule{\linewidth}{0.4pt}\vspace{2pt}
\end{abstract}

% ================================================================
\section{Introduction}
\label{sec:intro}

Take a fixed non-convex polygon, a 15-vertex shape that appeared in a
2024 Kaggle competition~\cite{kaggle2024}. Make $N$ identical copies.
The problem is to find the smallest square that holds all of them
without overlap. As $N$ grows, the configuration space expands
exponentially and exact enumeration becomes hopeless. Even for the
simpler case of unit disks the placement decision problem is
NP-complete~\cite{demaine2007}: if the decision version of a problem
is NP-complete then its optimization version is NP-hard, and ours is
at least as hard. Non-convexity adds a further layer of difficulty.
It makes collision detection more expensive, the energy landscape more
rugged, and the Minkowski sum characterizing pairwise overlap itself
non-convex and potentially disconnected.

Exact characterizations via the No-Fit Polygon (NFP), surveyed by
Bennell and Oliveira~\cite{bennell2008}, require convex decomposition
for non-convex shapes and become impractical at the scales we study.
Penalty-based SA consistently outperforms approaches restricted to
strictly feasible states: as Leung et al.~\cite{leung2003} showed on
orthogonal packing, hard feasibility constraints fragment the search
space into disconnected islands, while continuous penalties create
traversable bridges between them. Genetic algorithms can explore
structurally diverse configurations simultaneously~\cite{burke2006},
and Parallel Tempering~\cite{hukushima1996} mixes exponentially faster
than single-chain SA on multimodal landscapes, but both require
substantially longer budgets than we had available. As Sokal put it:
``Monte Carlo is an extremely bad method; it should be used only when
all alternative methods are worse.'' That is the situation here.

% ================================================================
\section{Problem Formulation}
\label{sec:formulation}

Let $\mathcal{P}$ be a fixed non-convex polygon with 15 vertices and
area $A_p = 0.245625$, computed by the shoelace formula from the
reference vertex array. Each copy $i$ is placed by specifying its
centre $(x_i, y_i)$ and rotation angle $\theta_i \in [0, 2\pi)$;
world vertices follow as $w_k^{(i)} = R(\theta_i)\,v_k + c_i$. A
configuration is \emph{feasible} if no two copies overlap and every
copy lies within $\Omega_L = [-L/2,\,L/2]^2$. The goal is to find,
for each $N$, the smallest $L$ admitting a feasible configuration.

Since all $N$ copies must fit inside $\Omega_L$ without overlap, the
square area satisfies $L^2 \geq N A_p$, giving the hard lower bound
$L_{\mathrm{area}} = \sqrt{N A_p}$; the bisection loop never probes
below this. Restricting the search to strictly non-overlapping
configurations at all times confines it to disconnected feasible
islands in $\mathbb{R}^{2N} \times [0, 2\pi)^N$. Allowing temporary,
penalized violations creates navigable paths between those islands,
which is why the choice of a continuous collision metric rather than a
binary one matters so much.

% ================================================================
\section{Algorithm}
\label{sec:algorithm}

\subsection{Collision Detection Pipeline}

Collision detection is invoked after every proposed move to evaluate
the energy change, making it the dominant runtime cost. Naive pairwise
checking costs $O(N^2)$ per iteration. A uniform spatial hash grid
with cell size $h = 2r_{\mathrm{br}} = 1.6$ (twice the bounding
radius $r_{\mathrm{br}} = 0.8$) restricts each query to a
$5 \times 5$ neighbourhood of ${\approx}\,261\eta$ candidates, where
$\eta = NA_p/L^2$ is the packing density. The energy change is then
computed incrementally, touching only pairs involving the moved
polygon, reducing the effective per-move cost to $O(1)$. Of those
candidates, roughly 96\% are eliminated by polygon-level
axis-aligned bounding box (AABB) checks; group-level and per-triangle
AABB checks prune the remainder before the Separating Axis Theorem
(SAT) is invoked.

The polygon is triangulated offline into $T = 13$ triangles, a fan
from vertex 0, grouped into three clusters of sizes $(5, 5, 3)$.
SAT projects a pair of triangles onto six edge-normal axes and
computes the minimum overlap width $d \geq 0$. If any axis shows a
gap the triangles are separated; when $d > 0$ it is the
\emph{penetration depth}, a continuous measure of overlap severity.
World-space vertex coordinates are projected directly onto each axis,
avoiding per-call recomputation of rotation matrices. A binary test
would assign equal cost to a barely touching pair and a deeply
embedded one, yielding a flat energy surface with no gradient to
follow. The depth $d$ provides a directional signal, and feeding
$d^2$ into the energy function creates the smooth guidance that steers
SA toward feasibility from arbitrary starting configurations.

\subsection{Energy Function and SA Schedule}

The total energy of a configuration at container side $L$ is
\begin{equation}
  E = \lambda \!\sum_{i<j}\sum_{(a,b)\in\mathcal{T}_{ij}}
  d(a,b)^{2}
  + \mu \sum_i \phi_i(L),
  \label{eq:energy}
\end{equation}
where $d(a,b)$ is the SAT penetration depth between triangle pair
$(a, b)$, $\phi_i(L)$ is the squared AABB boundary violation of copy
$i$, and $(\lambda, \mu)$ are penalty weights. Both terms are
quadratic in the geometric violation. A configuration is declared
feasible when the unweighted residual
$\mathcal{F} = \sum_{i<j} \text{overlap}_{ij} + \sum_i \phi_i < 10^{-6}$.

Each SA trial is time-budgeted rather than iteration-counted.
Temperature decreases geometrically,
$T_k = T_0\,\beta^k$ with $\beta = (T_{\mathrm{end}}/T_0)^{1/K}$,
from $T_0 = 1.0$ to $T_{\mathrm{end}} = 10^{-5}$ over $K = 100{,}000$
proposals. Step sizes $(s_{xy}, s_\theta)$ adapt every 2{,}000
iterations: if acceptance falls below 40\% they shrink by factor
0.95; above 60\% they grow by 1.05. This couples effective search
radius to temperature, producing large exploratory jumps at high $T$
and fine refinement near convergence. The penalty weights
$(\lambda, \mu)$ are held fixed within each trial and supplied by
the caller, which may ramp them between trials to progressively
tighten feasibility enforcement. Table~\ref{tab:sa} summarises the
parameters.

\begin{table}[t]
\centering\small
\setlength{\tabcolsep}{4.5pt}
\caption{SA schedule parameters (per trial).}
\label{tab:sa}
\begin{tabular}{lc}
\toprule
Parameter & Value \\
\midrule
Proposals per trial $K$   & $100{,}000$ \\
$T_{\mathrm{start}}$      & $1.0$ \\
$T_{\mathrm{end}}$        & $10^{-5}$ \\
Temperature schedule      & geometric decay \\
Step adapt window         & $2{,}000$ iterations \\
Acceptance target band    & $[0.40,\;0.60]$ \\
Step shrink / grow factor & $0.95$ / $1.05$ \\
Overlap penalty power $p$ & 2 (quadratic, fixed) \\
Feasibility threshold     & $\mathcal{F} < 10^{-6}$ \\
\bottomrule
\end{tabular}
\end{table}

\subsection{Outer Loop: Bisection and Polish}

The SA oracle answers one question: is there a feasible configuration
at this $L$? Finding the minimum $L$ requires an outer bisection loop.
The initial bracket is $[L_{\mathrm{lo}},\, L_{\mathrm{hi}}]$ where
$L_{\mathrm{lo}} = L_{\mathrm{area}} = \sqrt{N A_p}$ (the provably
infeasible lower bound) and $L_{\mathrm{hi}} = \gamma\,L_{\mathrm{lo}}$
with $\gamma \in \{2.2,\, 2.6,\, 3.0\}$ depending on $N$. At each
bisection step the SA is run at $L_{\mathrm{mid}} = \tfrac{1}{2}(
L_{\mathrm{lo}} + L_{\mathrm{hi}})$: if feasible, $L_{\mathrm{hi}}
\leftarrow L_{\mathrm{mid}}$; otherwise $L_{\mathrm{lo}} \leftarrow
L_{\mathrm{mid}}$. The loop terminates when the relative bracket width
satisfies $(L_{\mathrm{hi}} - L_{\mathrm{lo}}) \leq 10^{-3}\,
L_{\mathrm{hi}}$, certifying $L^*$ to a relative precision of
$10^{-3}$. Before each SA run at a new $L$, all polygon centres are
rescaled by $\alpha = L_{\mathrm{new}}/L_{\mathrm{old}}$, a warm
start that preserves structural information and reduces the time to
re-establish feasibility.

After bisection, a polish stage calls the SA oracle repeatedly at the
best known $L_{\mathrm{hi}}$ to improve configuration quality without
further narrowing the container. Progress is written to disk on every
improvement, and a global snapshot is flushed on \texttt{SIGTERM} so
SLURM preemption loses no work. Every worker is seeded
deterministically from (base seed, run ID, trial index) via
SplitMix64, guaranteeing full reproducibility.

% ================================================================
\section{Experimental Setup}
\label{sec:experiments}

All experiments ran on the University of Arizona HPC cluster (account
\texttt{ece569}, \texttt{standard} partition). Two complementary job
families were used.

The \emph{five-node embarrassingly parallel} jobs allocated five
exclusive nodes, each with 32 cores and 16~GB RAM. Worker count was
$W = 30$ per node (reserving two cores for the OS), giving 150
concurrent workers. Each node queued 64 seeds via \texttt{xargs -P 30},
so 320 seeds ran per job across three concurrency waves. This pipeline
produced the $N = 5$ and $N = 10$ cohorts (smoke tests, 3-minute wall
time) and was the intended design for larger $N$.

The \emph{single-node array} jobs ran at most two seeds concurrently
and produced the $N \in \{25, 50, 100, 200\}$ results reported here:
8 array tasks at $N = 25$ (4-hour wall time, effective budget
14{,}100\,s), 6 at $N = 50$ (4\,h), 4 at $N = 100$ (4\,h), and 2 at
$N = 200$ (8\,h, effective 28{,}500\,s). SLURM sent \texttt{SIGTERM}
120\,s before the hard kill, giving the solver time to flush its best
result. No inter-worker communication occurred during any run; the
global best for each $N$ was found in a post-hoc reduction pass over
all output CSVs.

% ================================================================
\section{Results}
\label{sec:results}

Table~\ref{tab:results} reports the best feasible side length $L^*$
and packing efficiency $\eta = N A_p / (L^*)^2$ for each $N$. The
area lower bound $L_{\mathrm{area}} = \sqrt{N A_p}$ is given for
reference; $\eta = 1$ would require zero wasted space, which is
geometrically impossible for this non-convex shape. The column
``Runs'' counts workers that found at least one feasible configuration.

\begin{table}[h]
\centering\small
\setlength{\tabcolsep}{3.5pt}
\caption{Best feasible side length, packing efficiency, and
feasibility rate by instance size.}
\label{tab:results}
\begin{tabular}{r@{\hskip6pt}ccccc}
\toprule
$N$ & $L_{\mathrm{area}}$ & $L^*$ & $\eta$ & Runs &
$t_{\mathrm{first}}$ \\
\midrule
  5 & 1.108 & 1.472 & 0.567 & 15 &  5\,s \\
 10 & 1.568 & 2.056 & 0.581 & 15 &  6\,s \\
 25 & 2.478 & 3.315 & 0.559 &  6 &  8\,s \\
 50 & 3.504 & 4.652 & 0.567 &  4 & 46\,s \\
100 & 4.956 & 6.785 & 0.533 &  3 & 53\,s \\
200 & 7.009 & 10.104& 0.481 &  1 & 19\,s \\
\bottomrule
\end{tabular}
\end{table}

Three patterns emerge. First, the number of workers reaching
feasibility drops sharply: 15 of 15 at $N = 5$ and $N = 10$, down to
6 at $N = 25$, and to a single run at $N = 200$. The SA is
time-budgeted and not iteration-bounded, but the cost per SA proposal
grows with $N$ because more polygon pairs must be checked. At large
$N$, fewer trials fit within the wall-time budget, and the trials that
do run explore a configuration space of dimension $3N$ with
proportionally less coverage per polygon. The configuration space at
$N = 200$ has 600 degrees of freedom; the budget was calibrated on
much smaller instances.

Second, best-of-best efficiency $\eta$ declines monotonically, from
0.58 at small $N$ to 0.48 at $N = 200$. Fitting $\eta(N) \approx
\eta_\infty + c/\!\sqrt{N}$ to the $N = 50$ and $N = 200$ points
gives $\eta_\infty \approx 0.40$ and $c \approx 1.22$. The scaling
has a clean physical basis: wasted space due to boundary proximity
grows with the container perimeter, which scales as $L \propto
\sqrt{N}$, while bulk packing capacity grows with the area, scaling
as $L^2 \propto N$; the relative boundary penalty therefore decays as
$1/L \propto 1/\!\sqrt{N}$. The extrapolated limit $\eta_\infty
\approx 0.40$ should be read with caution: at large $N$ the density
drop reflects solver budget limitations as much as any geometric
asymptote.

Third, the median time to final best configuration was 205\,s at
$N = 25$ but exceeded 3{,}250\,s at $N = 100$, out of a 14{,}100\,s
budget. Most improvement arrived at 23--26\% into the run, confirming
that the polish phase, not the bisection loop, produces the majority
of the density gain at large $N$. At small $N$ the bisection alone is
nearly sufficient; at large $N$ the bulk of the productive computation
happens during extended feasibility polish.
Figures~\ref{fig:timeseries}--\ref{fig:scaling} show these dynamics.

% ================================================================
\section{Conclusion}
\label{sec:conclusion}

Choosing SAT over a binary overlap test was the key design decision:
penetration depth provides the continuous energy gradient that steers
SA toward feasibility from arbitrary, heavily overlapping initial
configurations. The bisection loop tightens the bracket on $L^*$ to a
relative precision of $10^{-3}$, and the polish stage then fills
whatever wall time remains with incremental quality improvement. The
experiments confirm the solver reaches feasible packings for $N$ up to
200 within practical HPC allocations, but expose a clear bottleneck:
as $N$ grows, the time-budgeted SA visits fewer trials relative to the
dimension of the search space, and the feasibility rate collapses. The
natural first intervention is to scale the time budget proportionally
with $N$.

Three directions motivate future work. A CUDA implementation targeting
the P100 (compute capability 6.0) could exploit the embarrassingly
parallel trial structure, with one thread block per polygon pair and
triangle vertex arrays in shared memory. Parallel Tempering, which
runs replicas at a geometric temperature ladder and proposes Metropolis
swaps between adjacent temperatures, could seed cold-chain refinement
from configurations that single-chain SA never visits, though it
requires substantially longer wall times than tested here. Finally, a
genetic algorithm treating full polygon layouts as chromosomes could
maintain structural diversity across qualitatively different
configurations simultaneously, something serial SA cannot do.

% ================================================================
\bibliographystyle{plain}
\begin{thebibliography}{7}

\bibitem{kaggle2024}
Kaggle.
Santa 2024 competition: The polytope permutation puzzle.
\textit{kaggle.com/competitions/santa-2024}, 2024.

\bibitem{demaine2007}
E.~D. Demaine and J.~O'Rourke.
\textit{Geometric Folding Algorithms: Linkages, Origami, Polyhedra}.
Cambridge University Press, 2007.

\bibitem{bennell2008}
J.~A. Bennell and J.~F. Oliveira.
The geometry of nesting problems: A tutorial.
\textit{Eur.\ J.\ Oper.\ Res.}, 184(2):397--415, 2008.

\bibitem{leung2003}
T.~W. Leung, C.~K. Chan, and M.~D. Troutt.
Application of a mixed simulated annealing-genetic algorithm heuristic
for the two-dimensional orthogonal packing problem.
\textit{Eur.\ J.\ Oper.\ Res.}, 145(3):530--542, 2003.

\bibitem{burke2006}
E.~K. Burke, G.~Kendall, and G.~Whitwell.
A new placement heuristic for the orthogonal stock-cutting problem.
\textit{Oper.\ Res.}, 52(4):655--671, 2004.

\bibitem{hukushima1996}
K.~Hukushima and K.~Nemoto.
Exchange Monte Carlo method and application to spin glass simulations.
\textit{J.\ Phys.\ Soc.\ Jpn.}, 65(6):1604--1608, 1996.

\bibitem{kirkpatrick1983}
S.~Kirkpatrick, C.~D. Gelatt, and M.~P. Vecchi.
Optimization by simulated annealing.
\textit{Science}, 220(4598):671--680, 1983.

\end{thebibliography}

% ================================================================
% FIGURES: pages 4 and 5
% ================================================================
\clearpage

\begin{figure*}[p]
  \centering
  \includegraphics[width=\linewidth,height=0.88\textheight,keepaspectratio]%
    {2025/old/img/feasible_L_timeseries_N5_N10_N25_N50_N100_N200}
  \caption{Feasible square side length $L$ over wall time for $N \in
  \{5, 10, 25, 50, 100, 200\}$ (panels a--f). Thin lines are
  individual worker traces; the thick line is the aggregate
  best-so-far envelope. The right panel for each row shows the best
  packing geometry found. At $N = 200$ only one worker reached
  feasibility, producing a single visible trace.}
  \label{fig:timeseries}
\end{figure*}

\begin{figure*}[p]
  \centering
  \begin{minipage}[t]{0.555\linewidth}
    \includegraphics[width=\linewidth]%
      {2025/old/img/feasible_best_density_timeseries}
    \caption{Best-so-far packing efficiency $\eta = N A_p / L^2$
    versus wall time. Thin lines are individual workers; the thick
    line is the median across feasible workers. The plateau reflects
    the bisection loop locking in a bracket; improvement thereafter
    comes from the polish stage.}
    \label{fig:density}
  \end{minipage}
  \hfill
  \begin{minipage}[t]{0.425\linewidth}
    \includegraphics[width=\linewidth]%
      {2025/old/img/feasible_time_to_best_by_n}
    \caption{Wall time to final best configuration by $N$. Boxes are
    interquartile ranges; dots are individual workers. The sharp
    increase at $N \geq 50$ confirms that the polish stage, not
    bisection, produces most improvement at scale.}
    \label{fig:timetobest}
  \end{minipage}

  \vspace{10pt}

  \begin{minipage}[t]{0.49\linewidth}
    \includegraphics[width=\linewidth]{analysis/plots/density_vs_N}
    \caption{Best packing efficiency $\eta$ versus $N$. The
    monotonic decline reflects both boundary geometry effects and
    a time budget that covers fewer trials per degree of freedom
    as $N$ grows.}
    \label{fig:densityN}
  \end{minipage}
  \hfill
  \begin{minipage}[t]{0.49\linewidth}
    \includegraphics[width=\linewidth]{analysis/plots/density_vs_inv_sqrtN}
    \caption{Packing efficiency $\eta$ versus $1/\!\sqrt{N}$. The
    linear trend supports $\eta(N) \approx \eta_\infty +
    c/\!\sqrt{N}$ with $\eta_\infty \approx 0.40$, grounded in the
    container perimeter-to-area scaling argument.}
    \label{fig:scaling}
  \end{minipage}
\end{figure*}

\end{document}